% !TEX program = xelatex
\documentclass[11pt,a4paper]{article}
\usepackage{algo-preamble}

\title{\textbf{L02 \textperiodcentered{} Ασυμπτωτική Ανάλυση}}
\author{Algorithms Class Hub}
\date{\small Αλγόριθμοι και Πολυπλοκότητα (K17) \textperiodcentered{} ΕΚΠΑ}

\begin{document}
\maketitle

\begin{center}\itshape\small
Η γλώσσα με την οποία συγκρίνουμε αλγορίθμους: $O$, $\Theta$, $\Omega$, $o$,
$\omega$, η ιεραρχία πολυπλοκοτήτων, και η οριακή προσέγγιση με L'Hôpital.
\end{center}

\section{Η ερώτηση που θέλουμε να απαντήσουμε}

\begin{quoteblock}
\textbf{Πώς επηρεάζεται η απόδοση ενός αλγόριθμου όταν ο όγκος των δεδομένων που
λαμβάνει αυξάνεται;}
\end{quoteblock}

Στο πρώτο μάθημα είπαμε ότι η πολυπλοκότητα εκφράζεται ως συνάρτηση της διάστασης
$n$. Εδώ θα μάθουμε να \textit{συγκρίνουμε} τέτοιες συναρτήσεις χωρίς να μας
ενδιαφέρουν σταθεροί συντελεστές ή όροι χαμηλότερης τάξης --- γιατί όταν το $n$
μεγαλώνει, μόνο ο \textbf{κυρίαρχος όρος} μετράει.

\begin{intuition}
Σκέψου ότι κρίνεις δύο φοιτητές για τα μαθηματικά: ο ένας λύνει σύνθετα
προβλήματα σε 5 λεπτά, ο άλλος σε 5 λεπτά και 30 δευτερόλεπτα. Αν όλη η ζωή ήταν
μαθηματικά, η διαφορά θα ήταν ασήμαντη. Αλλά αν ο πρώτος λύνει σε 5 λεπτά και ο
δεύτερος σε \textbf{2 ώρες}, ο πρώτος είναι ασύγκριτα καλύτερος --- όχι κατά λίγο,
αλλά σε \textit{διαφορετική κλίμακα}. Η ασυμπτωτική ανάλυση μετρά ακριβώς αυτή τη
διαφορά κλίμακας.
\end{intuition}

\section{Γιατί δεν αρκούν οι σταθερές}

\begin{examplebox}
\textbf{Δύο αλγόριθμοι για το ίδιο πρόβλημα.}
\begin{itemize}
  \item $A_1$ τρέχει σε $T_1(n) = 100n$ στοιχειώδεις πράξεις.
  \item $A_2$ τρέχει σε $T_2(n) = n^2$.
\end{itemize}

Για ποιο $n$ είναι καλύτερος ο καθένας; Λύνουμε $100n = n^2 \Rightarrow n = 100$.
Άρα:
\begin{itemize}
  \item Για $n < 100$ ο $A_2$ είναι ταχύτερος.
  \item Για $n = 100$ είναι ισοδύναμοι.
  \item Για $n > 100$ ο $A_1$ είναι ταχύτερος --- και η διαφορά
  \textbf{μεγαλώνει συνεχώς}.
\end{itemize}

Αν θέσουμε ότι μία στοιχειώδης πράξη παίρνει $10^{-6}$ sec:

\begin{center}
\begin{tabular}{@{}l l l l@{}}
\toprule
$n$ & $A_1$ ($100n$) & $A_2$ ($n^2$) & Λόγος $A_2/A_1$ \\
\midrule
$1.000$   & $0.1$ sec & $1$ sec & $10\times$ \\
$100.000$ & $10$ sec  & $10.000$ sec $\approx 2.8$ ώρες & $1.000\times$ \\
\bottomrule
\end{tabular}
\end{center}

Παρά τη μεγάλη σταθερά ($100$) στον $A_1$, για μεγάλο $n$ ο $A_1$ νικάει σε
ακαριαία βάση. \textbf{Η σταθερά δεν σώζει έναν αλγόριθμο με χειρότερη τάξη
μεγέθους.}
\end{examplebox}

Αυτό το παράδειγμα είναι ο λόγος που η ασυμπτωτική ανάλυση \textbf{αγνοεί}
σταθερές. Αν η μία μέτρηση είναι $100n$ και η άλλη $n^2$, η σταθερά $100$ είναι
μικρή σε σχέση με τη διαφορά «γραμμικό vs τετραγωνικό».

\section{Οπτική σύγκριση συναρτήσεων}

Φαντάσου ότι σχεδιάζουμε στον ίδιο άξονα τις κλασικές κλάσεις πολυπλοκότητας ---
$\log n$, $n$, $n \log n$, $n^2$ και $2^n$. Παρατήρησε πώς η εκθετική $2^n$
ξεπερνά κάθε λογικό όριο για $n \approx 13$, ενώ η λογαριθμική κάμπτεται σχεδόν
αμέσως και μένει σχεδόν επίπεδη. Αυτή η οπτική διαφορά κλίμακας είναι όλο το
παιχνίδι.

\section{Οι τυπικοί ορισμοί}

Έστω $f, T : \mathbb{N} \to \mathbb{R}^+$. Ορίζουμε \textbf{τρεις κλάσεις
συναρτήσεων} σε σχέση με την $f$:

\[
O(f(n)) = \{T : \exists n_0 \in \mathbb{N},\ \exists c \in \mathbb{R}^+,\
\forall n \ge n_0,\ T(n) \le c \cdot f(n)\}
\]
\[
\Omega(f(n)) = \{T : \exists n_0 \in \mathbb{N},\ \exists c \in \mathbb{R}^+,\
\forall n \ge n_0,\ T(n) \ge c \cdot f(n)\}
\]
\[
\Theta(f(n)) = O(f(n)) \cap \Omega(f(n))
\]

Σε απλά Ελληνικά:
\begin{itemize}
  \item $T(n) = O(f(n))$ --- η $T$ «δεν αυξάνεται γρηγορότερα» από την $f$. Άνω
  φράγμα.
  \item $T(n) = \Omega(f(n))$ --- η $T$ «δεν αυξάνεται πιο αργά» από την $f$. Κάτω
  φράγμα.
  \item $T(n) = \Theta(f(n))$ --- η $T$ αυξάνεται \textbf{στον ίδιο ρυθμό} με την
  $f$ (μέχρι σταθεράς).
\end{itemize}

\begin{keypoint}
\textbf{Σταθερά $+$ $n_0$ είναι ο πυρήνας του ορισμού.} Η συνθήκη ισχύει
«τελικά» --- από κάποιο σημείο $n_0$ και μετά --- και «έως σταθεράς $c$». Έτσι
σταθερές και αρχικές διαφορές δεν μετράνε. Μετράει μόνο η \textbf{ασυμπτωτική}
συμπεριφορά.
\end{keypoint}

\section{Παραδείγματα ορισμού}

\begin{examplebox}
\textbf{Πρόβλημα.} Έστω $T(n) = 3n^2 - 100n + 6$. Δείξε ότι:
\begin{enumerate}
  \item $T(n) \in O(n^2)$ --- γιατί $3n^2 - 100n + 6 < 3n^2$ για κάθε $n \ge 1$.
  (Πάρε $c = 3$, $n_0 = 1$.)
  \item $T(n) \in O(n^3)$ --- για κάθε σταθερά $c > 0$, ισχύει
  $3n^2 - 100n + 6 < c \cdot n^3$ για αρκετά μεγάλο $n$. Δηλαδή το $O(n^3)$
  ισχύει αλλά είναι «ξεκάθαρα ασθενέστερο» από το $O(n^2)$.
  \item $T(n) \notin O(n)$ --- γιατί για κάθε σταθερά $c$, ισχύει $c \cdot n <
  3n^2$ για $n > c$. Άρα ο $T$ ξεπερνά τελικά κάθε γραμμική σταθερά.
\end{enumerate}

\textbf{Μάθημα:} όταν λέμε «$T(n) = O(g(n))$» δεν λέμε \textit{την ακριβή τάξη},
λέμε \textit{ένα άνω φράγμα}. Το $O(n^3)$ ισχύει αλλά είναι χαλαρό. Στο εξεταστικό
ζητάμε πάντα \textbf{το «καλύτερο» (πιο σφιχτό) $O$}, που συνήθως είναι το
$\Theta$.
\end{examplebox}

\section{Η ιεραρχία πολυπλοκοτήτων}

Σε αυξανόμενη σειρά ρυθμού αύξησης, οι κλασικές κλάσεις είναι:

\[
\underbrace{1 \prec \log n \prec n \prec n \log n \prec n^2}_{\text{πολυωνυμικό (καλό)}}
\prec
\underbrace{2^n \prec n!}_{\text{εκθετικό (κακό)}}
\]

Κάθε κλάση «κυριαρχείται» από την επόμενη: για κάθε διαδοχικό ζευγάρι, η αριστερή
ανήκει στο $o(\text{δεξιά})$.

\begin{keypoint}
\textbf{Ο πρακτικός κανόνας του μαθήματος:}
\begin{itemize}
  \item \textbf{Πολυωνυμική} πολυπλοκότητα ($n^c$ για σταθερά $c$) $\to$
  \textbf{καλή}.
  \item \textbf{Εκθετική} ($a^n$ για $a > 1$) $\to$ \textbf{κακή}.
\end{itemize}
Η διαχωριστική γραμμή είναι αυτή. Όλοι οι αλγόριθμοι που θα δούμε στοχεύουν
πολυωνυμική πολυπλοκότητα.
\end{keypoint}

\section{Wheat and chessboard --- γιατί το εκθετικό είναι «καταστροφή»}

\begin{examplebox}
\textbf{Ο μύθος.} Λέγεται ότι ο εφευρέτης του σκακιού ζήτησε ως αμοιβή από τον
βασιλιά: 1 κόκκο σιταριού στο πρώτο τετράγωνο, 2 στο δεύτερο, 4 στο τρίτο, \dots
διπλασιάζοντας μέχρι το 64ο τετράγωνο. Ο βασιλιάς, που νόμιζε ότι το αίτημα ήταν
ασήμαντο, δέχτηκε. Το συνολικό σιτάρι:

\[
T_{64} = 1 + 2 + 4 + \dots + 2^{63} = 2^{64} - 1 \approx 1.8 \times 10^{19}
\]

Δηλαδή \textbf{18 δισ.\ δισ.\ κόκκοι} --- περισσότερο σιτάρι από όσο έχει
παραχθεί ποτέ στην ιστορία της ανθρωπότητας. Ο βασιλιάς δεν είχε τρόπο να
πληρώσει.
\end{examplebox}

Αυτό είναι το $2^n$ σε δράση. Αν ένας αλγόριθμος είναι $O(2^n)$, αρκεί να αυξήσεις
το $n$ κατά 1 για να \textbf{διπλασιαστεί} ο χρόνος εκτέλεσης.

\section{Η «πικρή πραγματικότητα»: γιατί δεν μας σώζουν οι ταχύτεροι υπολογιστές}

Μια έντονη παρανόηση είναι ότι «οι αλγόριθμοι δεν έχουν τόση σημασία γιατί οι
υπολογιστές γίνονται όλο και ταχύτεροι». Ας το ελέγξουμε: \textbf{έστω υπολογιστής
1000 φορές ταχύτερος από τον σημερινό}. Πόσο μεγαλύτερο πρόβλημα μπορούμε να
λύσουμε στον ίδιο χρόνο;

\begin{center}
\begin{tabular}{@{}l l l@{}}
\toprule
Πολυπλοκότητα & Τωρινό μέγεθος & Μελλοντικό μέγεθος ($\times 1000$ ταχύτητα) \\
\midrule
$O(n)$   & $n$ & $1000n$ \\
$O(n^2)$ & $n$ & $n \cdot \sqrt{1000} \approx 31.6\,n$ \\
$O(n^3)$ & $n$ & $n \cdot 1000^{1/3} \approx 10n$ \\
$O(n^4)$ & $n$ & $n \cdot 1000^{1/4} \approx 5.6n$ \\
$O(2^n)$ & $n$ & $n + \log_2 1000 \approx n + 10$ \\
$O(3^n)$ & $n$ & $n + \log_3 1000 \approx n + 6.3$ \\
$O(n!)$  & $n$ & $\approx n$ (ουσιαστικά αμετάβλητο) \\
\bottomrule
\end{tabular}
\end{center}

\begin{warningbox}
Με $1000\times$ ταχύτερο μηχάνημα, ένας \textbf{γραμμικός} αλγόριθμος λύνει 1000
φορές μεγαλύτερο πρόβλημα. Ένας \textbf{εκθετικός} αλγόριθμος λύνει πρόβλημα
\textbf{10 παραπάνω εισόδων}. Καμία υλική ταχύτητα δεν σώζει έναν εκθετικό
αλγόριθμο. Γι' αυτό μετράει ο αλγόριθμος, όχι το μηχάνημα.
\end{warningbox}

\section{Ιδιότητες των τάξεων μεγέθους}

Έστω $f, f_1, f_2 : \mathbb{N} \to \mathbb{R}^+$. Όλα τα παρακάτω αποδεικνύονται
κατευθείαν από τον ορισμό:

\begin{enumerate}
  \item $f(n) \in O(f(n))$ --- κάθε συνάρτηση «κυριαρχείται από τον εαυτό της».
  \item $f(n) \in \Theta(f(n))$ --- και «είναι ίσης τάξης με τον εαυτό της».
  \item $O(c \cdot f(n)) = O(f(n))$ --- οι σταθεροί συντελεστές απορροφώνται.
  \item $O(f(n) + f(n)) = O(f(n))$ --- διπλασιάζοντας μια συνάρτηση δεν αλλάζει
  τάξη.
  \item $O(f_1(n) + f_2(n)) = O(\max\{f_1(n), f_2(n)\})$ --- σε άθροισμα μετράει ο
  κυρίαρχος όρος.
  \item $O(f_1(n)) \cdot O(f_2(n)) = O(f_1(n) \cdot f_2(n))$ --- οι τάξεις
  πολλαπλασιάζονται.
  \item $f(n) \in O(g(n)) \iff g(n) \in \Omega(f(n))$ --- $O$ και $\Omega$ είναι
  δυϊκά.
  \item $f(n) \in \Theta(g(n)) \iff f(n) \in O(g(n))$ ΚΑΙ $g(n) \in O(f(n))$ ---
  το $\Theta$ θέλει και τις δύο κατευθύνσεις.
  \item $f(n) \in \Theta(g(n)) \iff g(n) \in \Theta(f(n))$ --- το $\Theta$ είναι
  συμμετρικό.
  \item $f(n) \in \Theta(g(n))$ και $g(n) \in \Theta(h(n)) \Rightarrow f(n) \in
  \Theta(h(n))$ --- μεταβατικότητα.
\end{enumerate}

\begin{examplebox}
\textbf{Απόδειξη της 5ης ιδιότητας} ($O(f_1 + f_2) = O(\max\{f_1, f_2\})$).

Έστω $g(n) \in O(f_1(n))$ και $g(n) \in O(f_2(n))$. Από τον ορισμό:
\[
\exists c_1, n_1 : \forall n \ge n_1,\ g(n) \le c_1 f_1(n)
\]
\[
\exists c_2, n_2 : \forall n \ge n_2,\ g(n) \le c_2 f_2(n)
\]
Πάρε $c_0 = \max\{c_1, c_2\}$, $n_0 = \max\{n_1, n_2\}$. Τότε για κάθε
$n \ge n_0$:
\[
g(n) \le c_0 (f_1(n) + f_2(n)) \le 2 c_0 \max\{f_1(n), f_2(n)\}
\]
Άρα $g(n) \in O(\max\{f_1(n), f_2(n)\})$. $\blacksquare$
\end{examplebox}

\section{Είναι κάθε δύο συναρτήσεις ασυμπτωτικά συγκρίσιμες;}

\textbf{Όχι.} Υπάρχουν συναρτήσεις που η μία δεν είναι ούτε $O$ ούτε $\Omega$ της
άλλης.

\begin{examplebox}
\textbf{Αντιπαράδειγμα.} Έστω $f(n) = n$ και $g(n) = n^{1 + \sin(n)}$. Ο εκθέτης
$1 + \sin(n)$ ταλαντώνεται μεταξύ $0$ και $2$:
\begin{itemize}
  \item Όταν $\sin(n) \approx 1$, η $g(n) \approx n^2$ (αυξάνεται πολύ πιο
  γρήγορα από $f$).
  \item Όταν $\sin(n) \approx -1$, η $g(n) \approx n^0 = 1$ (αυξάνεται πολύ πιο
  αργά από $f$).
\end{itemize}
Άρα δεν μπορούμε να πούμε ούτε «η $g$ είναι $O(n)$» ούτε «η $g$ είναι
$\Omega(n)$». \textbf{Αυτές οι δύο συναρτήσεις είναι ασυμπτωτικά ασύγκριτες.}
\end{examplebox}

\section{Πρακτικές ασκήσεις}

\begin{examplebox}
\textbf{Άσκηση 1.} Να δειχθεί ότι αν $a > 1$ και $f(n) \in O(\log_a n)$, τότε
$f(n) \in O(\log_2 n)$. (Με άλλα λόγια: η βάση του λογάριθμου δεν παίζει ρόλο στο
$O$.)

\textbf{Λύση.} Από τον ορισμό: $f(n) \le c \log_a n$ για κάποια σταθερά $c$. Με
αλλαγή βάσης: $\log_a n = \dfrac{\log_2 n}{\log_2 a}$. Άρα:
\[
f(n) \le c \cdot \frac{\log_2 n}{\log_2 a} = c' \cdot \log_2 n
\]
όπου $c' = c / \log_2 a$ είναι σταθερά (γιατί $a > 1$ δίνει $\log_2 a > 0$). Άρα
$f(n) \in O(\log_2 n)$. $\blacksquare$

\textbf{Συμπέρασμα.} Γράφουμε απλά $O(\log n)$ --- δεν χρειάζεται να γράψουμε τη
βάση.
\end{examplebox}

\begin{examplebox}
\textbf{Άσκηση 2.} Να δειχθεί ότι $\log(n!) \in O(n \log n)$.

\textbf{Λύση.} Παρατήρησε:
\[
n! = 1 \cdot 2 \cdot 3 \cdots n \le n \cdot n \cdot n \cdots n = n^n
\]
Αφού η λογαριθμική συνάρτηση είναι αύξουσα:
\[
\log(n!) \le \log(n^n) = n \log n
\]
Άρα $\log(n!) \in O(n \log n)$ με σταθερά $c = 1$. $\blacksquare$

\textbf{Bonus.} Η αντίστροφη κατεύθυνση ισχύει επίσης ($\log(n!) \in \Omega(n
\log n)$), οπότε $\log(n!) \in \Theta(n \log n)$ --- αυτό θα το χρειαστούμε όταν
αποδείξουμε ότι η ταξινόμηση με συγκρίσεις απαιτεί τουλάχιστον $\Omega(n \log n)$
συγκρίσεις (κάτω φράγμα).
\end{examplebox}

\begin{examplebox}
\textbf{Άσκηση 3.} Βρες τον καλύτερο (πιο σφιχτό) $O$-συμβολισμό για κάθε
έκφραση:

\textbf{(α)} $2 \log n - 4n + 3n \log n$. Ο κυρίαρχος όρος είναι ο $3n \log n$
(κάθε άλλος είναι χαμηλότερης τάξης). Άρα $O(n \log n)$ --- στην πραγματικότητα
και $\Theta(n \log n)$.

\textbf{(β)} $2 + 4 + \dots + 2n = 2(1 + 2 + \dots + n) = 2 \cdot
\frac{n(n+1)}{2} = n^2 + n$. Άρα $\Theta(n^2)$.

\textbf{(γ)} $2 + 4 + 8 + \dots + 2^n = 2^{n+1} - 2$. Άρα $\Theta(2^n)$ --- η
ταυτότητα της γεωμετρικής σειράς δίνει ακριβές αποτέλεσμα.
\end{examplebox}

\begin{examplebox}
\textbf{Άσκηση 4.} Δείξε ότι:

\textbf{(α)} $\sum_{k=1}^n k \cdot n \in O(n^3)$. Πράγματι
$\sum_{k=1}^n k \cdot n = n \sum_{k=1}^n k = n \cdot \frac{n(n+1)}{2} =
\frac{n^2(n+1)}{2} \le n^3$. Άρα $O(n^3)$.

\textbf{(β)} Αν $f(n) = O(n)$, τότε $(f(n))^2 = O(n^2)$. Αν $f(n) \le c \cdot n$,
τότε $(f(n))^2 \le c^2 \cdot n^2$. Άρα $(f(n))^2 = O(n^2)$.

\textbf{(γ)} $3^n \notin O(2^n)$. Λόγος: $\dfrac{3^n}{2^n} = (3/2)^n \to \infty$.
Άρα δεν υπάρχει σταθερά $c$ με $3^n \le c \cdot 2^n$ για κάθε $n \ge n_0$ --- το
$3^n$ ξεπερνά τελικά κάθε σταθερό πολλαπλάσιο του $2^n$.
\end{examplebox}

\section{Εναλλακτική προσέγγιση: οριακές σχέσεις}

Αντί να δείξουμε τυπικά «υπάρχει $c$ και $n_0$», μπορούμε να χρησιμοποιήσουμε
\textbf{όρια}. Έστω $f, g : \mathbb{N} \to \mathbb{R}^+$:

\[
\lim_{n \to \infty} \frac{f(n)}{g(n)} = \alpha \ne 0 \;\Rightarrow\; f \in \Theta(g)
\]
\[
\lim_{n \to \infty} \frac{f(n)}{g(n)} = 0 \;\Rightarrow\; f \in O(g)\ \text{και}\
f \notin \Theta(g)
\]
\[
\lim_{n \to \infty} \frac{f(n)}{g(n)} = \infty \;\Rightarrow\; g \in O(f)\
\text{και}\ g \notin \Theta(f)
\]

Οι μεσαίες περιπτώσεις (όριο $0$ ή $\infty$) δίνουν τους «αυστηρά μικρότερα»
συμβολισμούς:
\begin{itemize}
  \item $\lim \to 0$ $\to$ γράφουμε $f(n) \in o(g(n))$ («μικρό-ο»). Διαβάζεται «η
  $f$ είναι \textbf{αυστηρά} μικρότερης τάξης από τη $g$».
  \item $\lim \to \infty$ $\to$ γράφουμε $f(n) \in \omega(g(n))$
  («μικρό-ωμέγα»). Διαβάζεται «η $f$ είναι \textbf{αυστηρά} μεγαλύτερης τάξης».
\end{itemize}

Πιο τυπικά:
\[
o(f(n)) = \{g : \forall c > 0,\ \exists n_0,\ \forall n \ge n_0,\ g(n) < c \cdot
f(n)\}
\]
\[
\omega(f(n)) = \{g : \forall c > 0,\ \exists n_0,\ \forall n \ge n_0,\ g(n) > c
\cdot f(n)\}
\]

\begin{keypoint}
\textbf{Η διαφορά $O$ vs $o$:}
\begin{itemize}
  \item $g \in O(f)$ $\to$ υπάρχει \textit{κάποιο} $c$ με $g(n) \le c f(n)$
  τελικά. Η $g$ ίσως είναι «του ίδιου ρυθμού».
  \item $g \in o(f)$ $\to$ για \textit{κάθε} $c$ μικρό, ισχύει $g(n) < c f(n)$
  τελικά. Η $g$ είναι \textbf{αυστηρά} πιο αργή.
\end{itemize}
Το $\Theta$ αποκλείεται με $o$. Το $\Theta$ επιτρέπεται με $O$.
\end{keypoint}

\section{Υπενθύμιση --- Κανόνας του L'Hôpital}

Πολλά από τα όρια που χρειαζόμαστε δίνουν απροσδιοριστία $\infty/\infty$ ή $0/0$.
Σε αυτές τις περιπτώσεις:

\begin{quoteblock}
Αν $\lim_{x \to \infty} f(x) = +\infty$ και $\lim_{x \to \infty} g(x) = +\infty$,
και οι $f, g$ είναι παραγωγίσιμες, τότε:
\[
\lim_{x \to \infty} \frac{f(x)}{g(x)} = \lim_{x \to \infty} \frac{f'(x)}{g'(x)}
\]
\end{quoteblock}

Όταν η απροσδιοριστία επιμένει, εφαρμόζουμε ξανά. Το χρησιμοποιούμε αμέσως
παρακάτω.

\section{Δύο σημαντικά αποτελέσματα για τους ρυθμούς αύξησης}

\subsection{Εκθετικό κυριαρχεί πάνω σε πολυωνυμικό}

\begin{quoteblock}
\textbf{Κάθε εκθετική συνάρτηση με βάση μεγαλύτερη του 1 αυξάνεται γρηγορότερα από
κάθε πολυωνυμική συνάρτηση.}
\end{quoteblock}

Τυπικά, για κάθε $a > 1$ και $b \in \mathbb{R}$:
\[
\lim_{n \to \infty} \frac{n^b}{a^n} = 0 \quad \Leftrightarrow \quad n^b \in o(a^n)
\]

\begin{examplebox}
\textbf{Απόδειξη.} Εφαρμόζουμε L'Hôpital επανειλημμένα:
\[
\lim_{x \to \infty} \frac{x^b}{a^x}
= \lim_{x \to \infty} \frac{b \cdot x^{b-1}}{a^x \ln a}
= \lim_{x \to \infty} \frac{b(b-1) x^{b-2}}{a^x \ln^2 a}
= \dots
= \lim_{x \to \infty} \frac{b!}{a^x \ln^b a} = 0
\]
Μετά από $\lceil b \rceil$ εφαρμογές του L'Hôpital, ο αριθμητής γίνεται σταθερά
και ο παρονομαστής εξακολουθεί να αυξάνεται εκθετικά --- το όριο είναι $0$.
$\blacksquare$
\end{examplebox}

\subsection{Πολυωνυμικό κυριαρχεί πάνω σε πολυλογαριθμικό}

Μια συνάρτηση $f$ είναι \textbf{πολυλογαριθμικά φραγμένη} αν υπάρχει σταθερά $c$
με $f(n) \le \log^c n$.

\begin{quoteblock}
\textbf{Κάθε πολυωνυμική συνάρτηση αυξάνεται γρηγορότερα από κάθε πολυλογαριθμική.}
\end{quoteblock}

Τυπικά, για $a > 0,\ b > 0$:
\[
\lim_{n \to \infty} \frac{\log^b n}{n^a} = 0 \quad \Leftrightarrow \quad \log^b n
\in o(n^a)
\]

\begin{examplebox}
\textbf{Απόδειξη (περίληψη).} Θέλουμε να δείξουμε ότι υπάρχει $n_0$ τέτοιο ώστε
για κάθε $n \ge n_0$:
\[
(\ln n)^b \le n^a \iff b \ln \ln n \le a \ln n \iff \ln \ln n \le \tfrac{a}{b}
\ln n
\]
Με L'Hôpital (χρησιμοποιώντας ότι $(\ln x)' = 1/x$ και τον κανόνα αλυσίδας):
\[
\lim_{x \to \infty} \frac{\ln \ln x}{\tfrac{a}{b} \ln x}
= \frac{b}{a} \lim_{x \to \infty} \frac{1/(x \ln x)}{1/x}
= \frac{b}{a} \lim_{x \to \infty} \frac{1}{\ln x} = 0
\]
Άρα ο λόγος τείνει στο $0$, που σημαίνει ότι ο παρονομαστής τελικά «νικάει».
$\blacksquare$
\end{examplebox}

\begin{keypoint}
\textbf{Συνδυάζοντας τα δύο αποτελέσματα:}
\[
\log^c n \prec n^\epsilon \prec n^c \prec a^n \prec n!
\]
για οποιεσδήποτε σταθερές $c > 0,\ \epsilon > 0,\ a > 1$. Αυτή είναι η «βασική
ιεραρχία» που χρησιμοποιούμε διαρκώς.
\end{keypoint}

\section{Η έννοια του βέλτιστου αλγορίθμου}

\begin{quoteblock}
Για ένα πρόβλημα $\Pi$ ζητάμε \textbf{τον καλύτερο} αλγόριθμο που το επιλύει ---
αυτόν με την καλύτερη πολυπλοκότητα.
\end{quoteblock}

Πιο τυπικά: αν \textit{κάθε} αλγόριθμος που επιλύει το $\Pi$ έχει ασυμπτωτική
πολυπλοκότητα $\Omega(f(n))$ --- αυτό είναι \textbf{κάτω φράγμα του προβλήματος}
--- τότε ένας αλγόριθμος με πολυπλοκότητα $O(f(n))$ είναι \textbf{(ασυμπτωτικά)
βέλτιστος}: ακουμπάει το θεωρητικό όριο.

\begin{intuition}
\textbf{Παράδειγμα.} Για την ταξινόμηση συγκρίσεων, αποδεικνύεται ότι κάθε
αλγόριθμος χρειάζεται $\Omega(n \log n)$ συγκρίσεις στη χείριστη περίπτωση (κάτω
φράγμα). Η συγχωνευτική ταξινόμηση τρέχει σε $O(n \log n)$ --- άρα είναι
\textbf{ασυμπτωτικά βέλτιστη}. Θα το αποδείξουμε στο επόμενο μάθημα, για το Divide
\& Conquer.
\end{intuition}

\section{Κλείσιμο: πολλαπλασιασμός τετραγωνικών πινάκων}

Ένα εξαιρετικό παράδειγμα του γιατί η ασυμπτωτική ανάλυση είναι \textbf{ζωντανό}
ερευνητικό αντικείμενο.

\begin{examplebox}
\textbf{Πρόβλημα.} Δοθέντων δύο $n \times n$ πινάκων $A = (a_{ij})$,
$B = (b_{ij})$, υπολόγισε τον πίνακα $C = A \cdot B$, όπου:
\[
c_{ij} = \sum_{k=1}^n a_{ik} \cdot b_{kj}
\]

\textbf{Naive αλγόριθμος.}
\begin{codebox}
\begin{verbatim}
Για i = 1 έως n:
  Για j = 1 έως n:
    c_ij <- 0
    Για k = 1 έως n:
      c_ij <- c_ij + a_ik * b_kj
\end{verbatim}
\end{codebox}

\textbf{Ανάλυση.} Τρεις φωλιασμένοι βρόχοι μέχρι το $n$ $\to$ $n^3$
πολλαπλασιασμοί $\to$ $\Theta(n^3)$.

\textbf{Κάτω φράγμα.} Χρειαζόμαστε να υπολογίσουμε $n^2$ στοιχεία του $C$. Άρα
κάθε αλγόριθμος είναι τουλάχιστον $\Omega(n^2)$. \textbf{Δεν γνωρίζουμε αλγόριθμο
που να επιτυγχάνει $O(n^2)$.}
\end{examplebox}

Το \textbf{χάσμα} $[\Omega(n^2),\ O(n^3)]$ είναι από τα πιο διάσημα ανοιχτά
προβλήματα. Η ιστορία του:

\begin{center}
\begin{tabular}{@{}l l l@{}}
\toprule
Έτος & Ερευνητές & Πολυπλοκότητα \\
\midrule
1969     & Strassen                              & $O(n^{2.81})$ \\
1986     & Coppersmith \& Winograd               & $O(n^{2.379})$ \\
2011     & V.~V.~Williams                        & $O(n^{2.373})$ \\
2023     & Williams, Xu, Xu, Zhou                & $O(n^{2.371552})$ \\
\textbf{2025} & \textbf{Alman, Duan, Williams, Xu, Xu, Zhou} & $O(n^{2.371339})$ \\
\bottomrule
\end{tabular}
\end{center}

Κάθε δεκαετία μικρο-βελτιώσεις στον εκθέτη. Δεν ξέρουμε αν το $\Omega(n^2)$ είναι
το πραγματικό κάτω φράγμα --- αλλά \textbf{πολλοί ερευνητές πιστεύουν ότι είναι}
και ότι κάποτε θα φτάσουμε εκεί.

\begin{recapbox}
\textbf{Τι κρατάμε από αυτή τη διάλεξη}
\begin{enumerate}
  \item \textbf{$O, \Omega, \Theta$}: άνω, κάτω, και ακριβές ασυμπτωτικό φράγμα.
  Ο ορισμός βασίζεται σε σταθερά $c$ και αρχικό $n_0$.
  \item \textbf{Σταθερές δεν μετράνε}: $O(100n) = O(n)$. Μόνο ο \textbf{κυρίαρχος
  όρος} μετράει.
  \item \textbf{Πολυωνυμικό vs εκθετικό}: η μεγάλη διαχωριστική γραμμή. Ταχύτερος
  υπολογιστής δεν σώζει εκθετικό αλγόριθμο.
  \item \textbf{Ιεραρχία}: $1 \prec \log n \prec n \prec n \log n \prec n^2 \prec
  \dots \prec 2^n \prec n!$.
  \item \textbf{Όρια ως εργαλείο}: $\lim f/g = 0 \to f \in o(g)$· $\lim = c \to
  f \in \Theta(g)$· $\lim = \infty \to f \in \omega(g)$.
  \item \textbf{L'Hôpital}: λύνει σχεδόν όλα τα όρια που χρειαζόμαστε.
  \item \textbf{Βέλτιστος αλγόριθμος}: αν το κάτω φράγμα του προβλήματος είναι
  $\Omega(f)$ και υπάρχει αλγόριθμος $O(f)$, είναι ασυμπτωτικά βέλτιστος.
\end{enumerate}
\end{recapbox}

\lecturefooter
\end{document}
